\section{P}

% packen (to pack)
\begin{verbentry}{
  style = {default},
  toc = {packen (to pack)},
  name = {packen},
  english = {to pack},
  type = {regular},
  auxiliary = {haben}
}


\vspace{5pt}
\hrule
\vspace{5pt}

\begin{tabular}{rl}
\multicolumn{2}{l}{\textbf{PRÄSENS}}\\
ich & \textbf{packe}\\
du & \textbf{packst}\\
er/sie/es & \textbf{packt}\\
wir & \textbf{packen}\\
ihr & \textbf{packt}\\
sie/Sie & \textbf{packen}\\
\end{tabular}
\hfill
\begin{tabular}{rl}
\multicolumn{2}{l}{\textbf{PRÄTERITUM}}\\
ich & \textbf{packte}\\
du & \textbf{packtest}\\
er/sie/es & \textbf{packte}\\
wir & \textbf{packten}\\
ihr & \textbf{packtet}\\
sie/Sie & \textbf{packten}\\
\end{tabular}
\hfill
\begin{tabular}{rl}
\multicolumn{2}{l}{\textbf{IMPERATIV}}\\
& \textbf{pack(e) (du)!}\\
& \textbf{packen wir!}\\
& \textbf{packt (ihr)!}\\
& \textbf{packen Sie!}\\
\end{tabular}

\vspace{10pt}

\begin{tabular}{rl}
\multicolumn{2}{l}{\textbf{PERFEKT}}\\
& haben + \textbf{gepackt}\\
\end{tabular}
\hfill
\begin{tabular}{rl}
\multicolumn{2}{l}{\textbf{FUTURE I}}\\
& werden + \textbf{packen}\\
\end{tabular}
\hfill
\begin{tabular}{rl}
\multicolumn{2}{l}{\textbf{PLUSQUAMPERFEKT}}\\
& hatten + \textbf{gepackt}\\
\end{tabular}

\vspace{10pt}

\begin{tabular}{lll}
\multicolumn{3}{l}{\textbf{MIT PRÄPOSITIONEN}}\\
& \textbf{packen in} + Akk & (to pack into)\\
& \textbf{packen auf} + Akk & (to load onto)\\
\end{tabular}
\hfill
\begin{tabular}{lll}
\multicolumn{3}{l}{\textbf{TRENNBARE VERBEN}}\\
& \textbf{ein$|$packen} & (to pack up)\\
& \textbf{aus$|$packen} & (to unpack)\\
& \textbf{an$|$packen} & (to tackle / get started)\\
& \textbf{mit$|$packen} & (to help pack)\\
\end{tabular}
\vspace{-5pt}
\end{verbentry}

% parken (to park)
\begin{verbentry}{
  style = {default},
  toc = {parken (to park)},
  name = {parken},
  english = {to park},
  type = {regular},
  auxiliary = {haben}
}
 
    
     \vspace{5pt}

    \hrule
    
    \vspace{5pt}

  \begin{tabular}{rl}
     \multicolumn{2}{l}{\textbf{PRÄSENS} }\\
    ich & \textbf{parke}\\
    du & \textbf{parkst}\\
    er/sie/es & \textbf{parkt}\\
    wir & \textbf{parken}\\
    ihr & \textbf{parkt}\\
    sie/Sie & \textbf{parken}\\
  \end{tabular}
  \hfill
     \begin{tabular}{rl}
    \multicolumn{2}{l}{\textbf{PRÄTERITUM} }\\
    ich & \textbf{parkte}\\
    du & \textbf{parktest}\\
    er/sie/es & \textbf{parkte}\\
    wir & \textbf{parkten}\\
    ihr & \textbf{parktet}\\
    sie/Sie & \textbf{parkten}\\
  \end{tabular}
  \hfill
  \begin{tabular}{rl}
     \multicolumn{2}{l}{\textbf{IMPERATIV} }\\
   & \textbf{park(e) (du)!}\\
   & \textbf{parken wir!}\\
   & \textbf{parket (ihr)!}\\
   & \textbf{parken Sie!}\\
   & \vspace{10pt}
  \end{tabular}

    \vspace{10pt}

 \begin{tabular}{rl}
     \multicolumn{2}{l}{\textbf{PERFEKT}}\\
  &  haben + \textbf{geparkt}\\
  \end{tabular}
\hfill
   \begin{tabular}{rl}
   
    \multicolumn{2}{l}{\textbf{FUTURE I} }\\
  &  werden + \textbf{parken}\\
 
  \end{tabular}
\hfill
   \begin{tabular}{rl}
      \multicolumn{2}{l}{\textbf{PLUSQUAMPERFEKT}}\\
   & hatten + \textbf{geparkt}\\
  \end{tabular}
  
  \vspace{10pt}

\begin{tabular}{lll}
\multicolumn{3}{l}{\textbf{MIT PRÄPOSITIONEN}}\\
& \textbf{---} & \\
\end{tabular}
\hfill
\begin{tabular}{lll}
\multicolumn{3}{l}{\textbf{TRENNBARE VERBEN}}\\
& \textbf{---} & \\
\end{tabular}
\end{verbentry}

% passen (to fit / to suit)
\begin{verbentry}{
  style = {dat},
  toc = {passen (to fit / to suit)},
  name = {passen},
  english = {to fit / to suit},
  type = {regular},
  auxiliary = {haben}
}
 
    
     \vspace{5pt}

    \hrule
    
    \vspace{5pt}

  \begin{tabular}{rl}
     \multicolumn{2}{l}{\textbf{PRÄSENS} }\\
    ich & \textbf{passe}\\
    du & \textbf{passt}\\
    er/sie/es & \textbf{passt}\\
    wir & \textbf{passen}\\
    ihr & \textbf{passt}\\
    sie/Sie & \textbf{passen}\\
  \end{tabular}
  \hfill
     \begin{tabular}{rl}
    \multicolumn{2}{l}{\textbf{PRÄTERITUM} }\\
    ich & \textbf{passte}\\
    du & \textbf{passtest}\\
    er/sie/es & \textbf{passte}\\
    wir & \textbf{passten}\\
    ihr & \textbf{passtet}\\
    sie/Sie & \textbf{passten}\\
  \end{tabular}
  \hfill
  \begin{tabular}{rl}
     \multicolumn{2}{l}{\textbf{IMPERATIV} }\\
   & \textbf{pass(e) (du)!}\\
   & \textbf{passen wir!}\\
   & \textbf{passt (ihr)!}\\
   & \textbf{passen Sie!}\\
   & \vspace{10pt}
  \end{tabular}

    \vspace{10pt}

 \begin{tabular}{rl}
     \multicolumn{2}{l}{\textbf{PERFEKT}}\\
   & haben + \textbf{gepasst}\\
  \end{tabular}
\hfill
   \begin{tabular}{rl}
   
    \multicolumn{2}{l}{\textbf{FUTURE I} }\\
   & werden + \textbf{passen}\\
 
  \end{tabular}
\hfill
   \begin{tabular}{rl}
      \multicolumn{2}{l}{\textbf{PLUSQUAMPERFEKT}}\\
   & hatten + \textbf{gepasst}\\
  \end{tabular}

   \vspace{10pt}

   \begin{tabular}{lll}
      \multicolumn{3}{l}{\textbf{MIT PRÄPOSITIONEN}}\\
& \textbf{passen zu} + Dat & (to suit / go with)\\
& \textbf{passen in} + Akk & (to fit into)\\
& \textbf{passen auf} + Akk & (to watch / take care of)
\end{tabular}
  \hfill
   \begin{tabular}{lll}
      \multicolumn{3}{l}{\textbf{TRENNBARE VERBEN}}\\
 & \textbf{auf$|$passen} & (to pay attention / watch)\\
 & \textbf{an$|$passen} & (to adapt / adjust)\\
  \end{tabular}


  \vspace{-5pt}
\end{verbentry}

% passieren (to happen)
\begin{verbentry}{
  style = {default},
  toc = {passieren (to happen)},
  name = {passieren},
  english = {to happen},
  type = {irregular},
  auxiliary = {sein}
}
 
    
     \vspace{5pt}

    \hrule
    
    \vspace{5pt}

  \begin{tabular}{rl}
     \multicolumn{2}{l}{\textbf{PRÄSENS} }\\
    ich & \textbf{-}\\
    du & \textbf{-}\\
    er/sie/es & \textbf{passiert}\\
    wir & \textbf{-}\\
    ihr & \textbf{-}\\
    sie/Sie & \textbf{passieren}\\
  \end{tabular}
  \hfill
     \begin{tabular}{rl}
    \multicolumn{2}{l}{\textbf{PRÄTERITUM} }\\
    ich & \textbf{-}\\
    du & \textbf{-}\\
    er/sie/es & \textbf{passierte}\\
    wir & \textbf{-}\\
    ihr & \textbf{-}\\
    sie/Sie & \textbf{passierten}\\
  \end{tabular}
  \hfill
  \begin{tabular}{rl}
     \multicolumn{2}{l}{\textbf{IMPERATIV} }\\
   & \textbf{-}\\
   & \textbf{-}\\
   & \textbf{-}\\
   & \textbf{-}\\
   & \vspace{10pt}
  \end{tabular}

    \vspace{10pt}

 \begin{tabular}{rl}
     \multicolumn{2}{l}{\textbf{PERFEKT}}\\
   & sein + \textbf{passiert}\\
  \end{tabular}
\hfill
   \begin{tabular}{rl}
   
    \multicolumn{2}{l}{\textbf{FUTURE I} }\\
   & werden + \textbf{passieren}\\
 
  \end{tabular}
\hfill
   \begin{tabular}{rl}
      \multicolumn{2}{l}{\textbf{PLUSQUAMPERFEKT}}\\
   & waren + \textbf{passiert}\\
  \end{tabular}
  
  \vspace{10pt}



\begin{tabular}{lll}
\multicolumn{3}{l}{\textbf{MIT PRÄPOSITIONEN}}\\
& \textbf{---} & \\
\end{tabular}
\hfill
\begin{tabular}{lll}
\multicolumn{3}{l}{\textbf{TRENNBARE VERBEN}}\\
& \textbf{---} & \\
\end{tabular}

\vspace{-5pt}
\end{verbentry}

% picknicken (to have a picnic)
\begin{verbentry}{
  style = {default},
  toc = {picknicken (to have a picnic)},
  name = {picknicken},
  english = {to have a picnic},
  type = {regular},
  auxiliary = {haben}
}


\vspace{5pt}
\hrule
\vspace{5pt}

\begin{tabular}{rl}
\multicolumn{2}{l}{\textbf{PRÄSENS}}\\
ich & \textbf{picknicke}\\
du & \textbf{picknickst}\\
er/sie/es & \textbf{picknickt}\\
wir & \textbf{picknicken}\\
ihr & \textbf{picknickt}\\
sie/Sie & \textbf{picknicken}\\
\end{tabular}
\hfill
\begin{tabular}{rl}
\multicolumn{2}{l}{\textbf{PRÄTERITUM}}\\
ich & \textbf{picknickte}\\
du & \textbf{picknicktest}\\
er/sie/es & \textbf{picknickte}\\
wir & \textbf{picknickten}\\
ihr & \textbf{picknicktet}\\
sie/Sie & \textbf{picknickten}\\
\end{tabular}
\hfill
\begin{tabular}{rl}
\multicolumn{2}{l}{\textbf{IMPERATIV}}\\
& \textbf{picknicke (du)!}\\
& \textbf{picknicken wir!}\\
& \textbf{picknickt (ihr)!}\\
& \textbf{picknicken Sie!}\\
\end{tabular}

\vspace{10pt}

\begin{tabular}{rl}
\multicolumn{2}{l}{\textbf{PERFEKT}}\\
& haben + \textbf{gepicknickt}\\
\end{tabular}
\hfill
\begin{tabular}{rl}
\multicolumn{2}{l}{\textbf{FUTURE I}}\\
& werden + \textbf{picknicken}\\
\end{tabular}
\hfill
\begin{tabular}{rl}
\multicolumn{2}{l}{\textbf{PLUSQUAMPERFEKT}}\\
& hatten + \textbf{gepicknickt}\\
\end{tabular}

\vspace{10pt}

\begin{tabular}{lll}
\multicolumn{3}{l}{\textbf{MIT PRÄPOSITIONEN}}\\
& \textbf{picknicken im} + Dat & (to have a picnic in)\\
& \textbf{picknicken auf} + Dat & (to have a picnic on)\\
\end{tabular}
\hfill
\begin{tabular}{lll}
\multicolumn{3}{l}{\textbf{TRENNBARE VERBEN}}\\
& \textbf{---} & \\
\end{tabular}

\vspace{-5pt}
\end{verbentry}

% planen (to plan)
\begin{verbentry}{
  style = {acc},
  toc = {planen (to plan)},
  name = {planen},
  english = {to plan},
  type = {regular, + Akkusativ},
  auxiliary = {haben}
}
 
    
     \vspace{5pt}

    \hrule
    
    \vspace{5pt}

  \begin{tabular}{rl}
     \multicolumn{2}{l}{\textbf{PRÄSENS} }\\
    ich & \textbf{plane}\\
    du & \textbf{planst}\\
    er/sie/es & \textbf{plant}\\
    wir & \textbf{planen}\\
    ihr & \textbf{plant}\\
    sie/Sie & \textbf{planen}\\
  \end{tabular}
  \hfill
     \begin{tabular}{rl}
    \multicolumn{2}{l}{\textbf{PRÄTERITUM} }\\
    ich & \textbf{plante}\\
    du & \textbf{plantest}\\
    er/sie/es & \textbf{plante}\\
    wir & \textbf{planten}\\
    ihr & \textbf{plantet}\\
    sie/Sie & \textbf{planten}\\
  \end{tabular}
  \hfill
  \begin{tabular}{rl}
     \multicolumn{2}{l}{\textbf{IMPERATIV} }\\
   & \textbf{plan(e) (du)!}\\
   & \textbf{planen wir!}\\
   & \textbf{plant (ihr)!}\\
   & \textbf{planen Sie!}\\
   & \vspace{10pt}
  \end{tabular}

    \vspace{10pt}

 \begin{tabular}{rl}
     \multicolumn{2}{l}{\textbf{PERFEKT}}\\
   & haben + \textbf{geplant}\\
  \end{tabular}
\hfill
   \begin{tabular}{rl}
   
    \multicolumn{2}{l}{\textbf{FUTURE I} }\\
   & werden + \textbf{planen}\\
 
  \end{tabular}
\hfill
   \begin{tabular}{rl}
      \multicolumn{2}{l}{\textbf{PLUSQUAMPERFEKT}}\\
   & hatten + \textbf{geplant}\\
  \end{tabular}
  
  
\vspace{10pt}

\begin{tabular}{lll}
\multicolumn{3}{l}{\textbf{MIT PRÄPOSITIONEN}}\\
& \textbf{---} & \\
\end{tabular}
\hfill
\begin{tabular}{lll}
\multicolumn{3}{l}{\textbf{TRENNBARE VERBEN}}\\
& \textbf{---} & \\
\end{tabular}
\vspace{-5pt}
\end{verbentry}

% plaudern (to chat)
\begin{verbentry}{
  style = {default},
  toc = {plaudern (to chat)},
  name = {plaudern},
  english = {to chat},
  type = {regular},
  auxiliary = {haben}
}


\vspace{5pt}
\hrule
\vspace{5pt}

\begin{tabular}{rl}
\multicolumn{2}{l}{\textbf{PRÄSENS}}\\
ich & \textbf{plaudere}\\
du & \textbf{plauderst}\\
er/sie/es & \textbf{plaudert}\\
wir & \textbf{plaudern}\\
ihr & \textbf{plaudert}\\
sie/Sie & \textbf{plaudern}\\
\end{tabular}
\hfill
\begin{tabular}{rl}
\multicolumn{2}{l}{\textbf{PRÄTERITUM}}\\
ich & \textbf{plauderte}\\
du & \textbf{plaudertest}\\
er/sie/es & \textbf{plauderte}\\
wir & \textbf{plauderten}\\
ihr & \textbf{plaudertet}\\
sie/Sie & \textbf{plauderten}\\
\end{tabular}
\hfill
\begin{tabular}{rl}
\multicolumn{2}{l}{\textbf{IMPERATIV}}\\
& \textbf{plaudere (du)!}\\
& \textbf{plaudern wir!}\\
& \textbf{plaudert (ihr)!}\\
& \textbf{plaudern Sie!}\\
\end{tabular}

\vspace{10pt}

\begin{tabular}{rl}
\multicolumn{2}{l}{\textbf{PERFEKT}}\\
& haben + \textbf{geplaudert}\\
\end{tabular}
\hfill
\begin{tabular}{rl}
\multicolumn{2}{l}{\textbf{FUTURE I}}\\
& werden + \textbf{plaudern}\\
\end{tabular}
\hfill
\begin{tabular}{rl}
\multicolumn{2}{l}{\textbf{PLUSQUAMPERFEKT}}\\
& hatten + \textbf{geplaudert}\\
\end{tabular}

\vspace{10pt}

\begin{tabular}{lll}
\multicolumn{3}{l}{\textbf{MIT PRÄPOSITIONEN}}\\
& \textbf{plaudern mit} + Dat & (to chat with)\\
& \textbf{plaudern über} + Akk & (to chat about)\\
\end{tabular}
\hfill
\begin{tabular}{lll}
\multicolumn{3}{l}{\textbf{TRENNBARE VERBEN}}\\
& aus$|$plaudern & (to let slip / blurt out)\\
\vspace{10pt}
\end{tabular}

\vspace{-5pt}
\end{verbentry}

% posten (to post)
\begin{verbentry}{
  style = {default},
  toc = {posten (to post)},
  name = {posten},
  english = {to post},
  type = {regular},
  auxiliary = {haben}
}
 

\vspace{5pt}
\hrule
\vspace{5pt}

\begin{tabular}{rl}
\multicolumn{2}{l}{\textbf{PRÄSENS}}\\
ich & \textbf{poste}\\
du & \textbf{postest}\\
er/sie/es & \textbf{postet}\\
wir & \textbf{posten}\\
ihr & \textbf{postet}\\
sie/Sie & \textbf{posten}\\
\end{tabular}
\hfill
\begin{tabular}{rl}
\multicolumn{2}{l}{\textbf{PRÄTERITUM}}\\
ich & \textbf{postete}\\
du & \textbf{postetest}\\
er/sie/es & \textbf{postete}\\
wir & \textbf{posteten}\\
ihr & \textbf{postetet}\\
sie/Sie & \textbf{posteten}\\
\end{tabular}
\hfill
\begin{tabular}{rl}
\multicolumn{2}{l}{\textbf{IMPERATIV}}\\
& \textbf{poste (du)!}\\
& \textbf{posten wir!}\\
& \textbf{postet (ihr)!}\\
& \textbf{posten Sie!}\\
\end{tabular}

\vspace{10pt}

\begin{tabular}{rl}
\multicolumn{2}{l}{\textbf{PERFEKT}}\\
& haben + \textbf{gepostet}\\
\end{tabular}
\hfill
\begin{tabular}{rl}
\multicolumn{2}{l}{\textbf{FUTURE I}}\\
& werden + \textbf{posten}\\
\end{tabular}
\hfill
\begin{tabular}{rl}
\multicolumn{2}{l}{\textbf{PLUSQUAMPERFEKT}}\\
& hatten + \textbf{gepostet}\\
\end{tabular}

\vspace{10pt}

\begin{tabular}{lll}
\multicolumn{3}{l}{\textbf{MIT PRÄPOSITIONEN}}\\
& \textbf{posten auf} + Dat & (to post on)\\
\end{tabular}
\hfill
\begin{tabular}{lll}
\multicolumn{3}{l}{\textbf{TRENNBARE VERBEN}}\\
& \textbf{---} & \\
\end{tabular}

\vspace{-5pt}
\end{verbentry}

% probieren (to try)
\begin{verbentry}{
  style = {acc},
  toc = {probieren (to try)},
  name = {probieren},
  english = {to try},
  type = {irregular, + Akkusativ},
  auxiliary = {haben}
}
 
    
     \vspace{5pt}

    \hrule
    
    \vspace{5pt}

  \begin{tabular}{rl}
     \multicolumn{2}{l}{\textbf{PRÄSENS} }\\
    ich & \textbf{probiere}\\
    du & \textbf{probierst}\\
    er/sie/es & \textbf{probiert}\\
    wir & \textbf{probieren}\\
    ihr & \textbf{probiert}\\
    sie/Sie & \textbf{probieren}\\
  \end{tabular}
  \hfill
     \begin{tabular}{rl}
    \multicolumn{2}{l}{\textbf{PRÄTERITUM} }\\
    ich & \textbf{probierte}\\
    du & \textbf{probiertest}\\
    er/sie/es & \textbf{probierte}\\
    wir & \textbf{probierten}\\
    ihr & \textbf{probiertet}\\
    sie/Sie & \textbf{probierten}\\
  \end{tabular}
  \hfill
  \begin{tabular}{rl}
     \multicolumn{2}{l}{\textbf{IMPERATIV} }\\
   & \textbf{probier(e) (du)!}\\
   & \textbf{probieren wir!}\\
   & \textbf{probiert (ihr)!}\\
   & \textbf{probieren Sie!}\\
   & \vspace{10pt}
  \end{tabular}

    \vspace{10pt}

 \begin{tabular}{rl}
     \multicolumn{2}{l}{\textbf{PERFEKT}}\\
   & haben + \textbf{probiert}\\
  \end{tabular}
\hfill
   \begin{tabular}{rl}
   
    \multicolumn{2}{l}{\textbf{FUTURE I} }\\
   & werden + \textbf{probieren}\\
 
  \end{tabular}
\hfill
   \begin{tabular}{rl}
      \multicolumn{2}{l}{\textbf{PLUSQUAMPERFEKT}}\\
   & hatten + \textbf{probiert}\\
  \end{tabular}
  
  
\vspace{10pt}

\begin{tabular}{lll}
\multicolumn{3}{l}{\textbf{MIT PRÄPOSITIONEN}}\\
& \textbf{---} & \\
\end{tabular}
\hfill
\begin{tabular}{lll}
\multicolumn{3}{l}{\textbf{TRENNBARE VERBEN}}\\
& \textbf{---} & \\
\end{tabular}
\vspace{-5pt}
\end{verbentry}

% prüfen (to check / examine)
\begin{verbentry}{
  style = {default},
  toc = {prüfen (to check / examine)},
  name = {prüfen},
  english = {to check / examine},
  type = {regular},
  auxiliary = {haben}
}


\vspace{5pt}
\hrule
\vspace{5pt}

\begin{tabular}{rl}
\multicolumn{2}{l}{\textbf{PRÄSENS}}\\
ich & \textbf{prüfe}\\
du & \textbf{prüfst}\\
er/sie/es & \textbf{prüft}\\
wir & \textbf{prüfen}\\
ihr & \textbf{prüft}\\
sie/Sie & \textbf{prüfen}\\
\end{tabular}
\hfill
\begin{tabular}{rl}
\multicolumn{2}{l}{\textbf{PRÄTERITUM}}\\
ich & \textbf{prüfte}\\
du & \textbf{prüftest}\\
er/sie/es & \textbf{prüfte}\\
wir & \textbf{prüften}\\
ihr & \textbf{prüftet}\\
sie/Sie & \textbf{prüften}\\
\end{tabular}
\hfill
\begin{tabular}{rl}
\multicolumn{2}{l}{\textbf{IMPERATIV}}\\
& \textbf{prüfe (du)!}\\
& \textbf{prüfen wir!}\\
& \textbf{prüft (ihr)!}\\
& \textbf{prüfen Sie!}\\
\end{tabular}

\vspace{10pt}

\begin{tabular}{rl}
\multicolumn{2}{l}{\textbf{PERFEKT}}\\
& haben + \textbf{geprüft}\\
\end{tabular}
\hfill
\begin{tabular}{rl}
\multicolumn{2}{l}{\textbf{FUTURE I}}\\
& werden + \textbf{prüfen}\\
\end{tabular}
\hfill
\begin{tabular}{rl}
\multicolumn{2}{l}{\textbf{PLUSQUAMPERFEKT}}\\
& hatten + \textbf{geprüft}\\
\end{tabular}

\vspace{10pt}

\begin{tabular}{lll}
\multicolumn{3}{l}{\textbf{MIT PRÄPOSITIONEN}}\\
& \textbf{prüfen auf} + Akk & (to check for)\\
\end{tabular}
\hfill
\begin{tabular}{lll}
\multicolumn{3}{l}{\textbf{TRENNBARE VERBEN}}\\
& \textbf{---} & \\
\end{tabular}
\vspace{-5pt}
\end{verbentry}

% putzen (to clean)
\begin{verbentry}{
  style = {acc},
  toc = {putzen (to clean)},
  name = {putzen},
  english = {to clean},
  type = {regular, + Akkusativ},
  auxiliary = {haben}
}
 
    
     \vspace{5pt}

    \hrule
    
    \vspace{5pt}

  \begin{tabular}{rl}
     \multicolumn{2}{l}{\textbf{PRÄSENS} }\\
    ich & \textbf{putze}\\
    du & \textbf{putzt}\\
    er/sie/es & \textbf{putzt}\\
    wir & \textbf{putzen}\\
    ihr & \textbf{putzt}\\
    sie/Sie & \textbf{putzen}\\
  \end{tabular}
  \hfill
     \begin{tabular}{rl}
    \multicolumn{2}{l}{\textbf{PRÄTERITUM} }\\
    ich & \textbf{putzte}\\
    du & \textbf{putztest}\\
    er/sie/es & \textbf{putzte}\\
    wir & \textbf{putzten}\\
    ihr & \textbf{putztet}\\
    sie/Sie & \textbf{putzten}\\
  \end{tabular}
  \hfill
  \begin{tabular}{rl}
     \multicolumn{2}{l}{\textbf{IMPERATIV} }\\
   & \textbf{putz(e) (du)!}\\
   & \textbf{putzen wir!}\\
   & \textbf{putzt (ihr)!}\\
   & \textbf{putzen Sie!}\\
   & \vspace{10pt}
  \end{tabular}

    \vspace{10pt}

 \begin{tabular}{rl}
     \multicolumn{2}{l}{\textbf{PERFEKT}}\\
 &  haben + \textbf{geputzt}\\
  \end{tabular}
\hfill
   \begin{tabular}{rl}
   
    \multicolumn{2}{l}{\textbf{FUTURE I} }\\
  &  werden + \textbf{putzen}\\
 
  \end{tabular}
\hfill
   \begin{tabular}{rl}
      \multicolumn{2}{l}{\textbf{PLUSQUAMPERFEKT}}\\
   & hatten + \textbf{geputzt}\\
  \end{tabular}
  
  
\vspace{10pt}

\begin{tabular}{lll}
\multicolumn{3}{l}{\textbf{MIT PRÄPOSITIONEN}}\\
& \textbf{---} & \\
\end{tabular}
\hfill
\begin{tabular}{lll}
\multicolumn{3}{l}{\textbf{TRENNBARE VERBEN}}\\
& \textbf{---} & \\
\end{tabular}
\vspace{-5pt}
\end{verbentry}
